% Special D=210 cocycle audit. This does not disprove the desired height bound.
\section{The compatible degree character and scalar splitting obstruction}

We retain the established moduli and tame-branch description of
$X_0(210,1)\to X_0(210,1)/W$: the branch at $s=0$ has inertia
$w_{105}$. The level-three cover defines all endomorphisms and compatible
isogenies over its fiber splitting field. The purpose here is to determine
the degree character of those compatible isogenies, and then to test the
proposed uniformly bounded scalar splitting coefficient field.

\subsection{The degree character is the Atkin--Lehner norm character}
Use the published complex uniformization of \cite[Chapter 43,
equations (43.1.2)--(43.1.3), Main Theorem 43.6.14]{Voight2021}:
\[
 v_z=(z,1)^t,\qquad\Lambda_z=\mathcal O v_z,\qquad
 A_z=\mathbb C^2/\Lambda_z.
\]
The left action of $\mathcal O$ on $\mathbb C^2$ is its specified QM
action. Let $b\in N_{B^\times}(\mathcal O)$ have positive norm and
write $b=\left(\begin{smallmatrix}a&b'\\c&d\end{smallmatrix}\right)$
under the real matrix embedding. Then
\[
 b v_z=(cz+d)v_{bz},\qquad
 (cz+d)\Lambda_{bz}=\mathcal O b v_z.
\]
Scale $b$ by a positive rational number if necessary so that
$\mathcal O b\subset\mathcal O$. The scalar complex-linear map
\[
 (cz+d)I_2:A_{bz}\longrightarrow A_z
\]
is an isogeny and commutes with the whole specified $\mathcal O$ action.
Its degree is
\[
 [\mathcal O:\mathcal O b]=\operatorname{Nrd}(b)^2.
\]
Indeed, right multiplication by $b$ on the four-dimensional rational
vector space $B$ has determinant $\operatorname{Nrd}(b)^2$; its image
lattice is $\mathcal O b$. Thus this is the fully endomorphism-compatible
isogeny, not just an unmarked isomorphism of abelian surfaces.

For comparison, the complex-linear map $(cz+d)b^{-1}$ identifies the
underlying lattices because $b^{-1}\mathcal O b=\mathcal O$, but twists
the QM action by an inner automorphism. Using that unmarked map would
give the wrong degree character. The scalar map is the required one.

In the non-CM case $\operatorname{End}^0(A_z)=B$. Any two compatible
quasi-isogenies differ by an element of the centralizer of $B$ in $B$,
namely a rational scalar. Their positive square-root degrees therefore
have the same class in $\mathbb Q^\times/\mathbb Q^{\times2}$.
For the coset $w_m$, the normalizer norm has square class $m$. The
level-change kernel changes only a marking and has square class one.
Consequently the intrinsic degree character sends $w_m$ to $m$.

\subsection{The sign class at an odd branch contact}
Let $\mu_\sigma:{}^\sigma A\to A$ be compatible quasi-isogenies and
let $c(\sigma,\tau)\in\mathbb Q^\times$ be their scalar cocycle.
Put $d_\sigma=\sqrt{\deg\mu_\sigma}>0$. These numbers are rational
by the calculation just made, and
\[
 c(\sigma,\tau)^2=d_\sigma d_\tau/d_{\sigma\tau}.
\]
Write its square-class character as
$d_\sigma=\mathbb Q^{\times2}\prod_jm_j^{\chi_j(\sigma)}$, with
$m_j\in\{2,3,5,7\}$ and $\chi_j$ quadratic, and write $a_j$ for
the corresponding quadratic extensions. The differential descent
calculation in the accompanying descent audit gives
\[
 [\operatorname{sign}c]=[B]+\sum_j(m_j,a_j)\quad\hbox{in }\operatorname{Br}(\mathbb Q)[2].
\]
For clarity, its two ingredients are as follows. Differentials of the
$\mu_\sigma$ give a semilinear descent datum on $M_2(L)$ whose invariant
algebra is the given copy of $B$. This identifies the natural-action
Brauer image of $c$ with $[B]$. Separately, the trivial-action coboundary
of the positive square root of $d_\sigma$ has natural-action image
$\sum_j(m_j,a_j)$. Dividing by that positive cocycle leaves the displayed
sign class.

Let $p>76176$ be an odd prime outside the fixed bad set of the tame cover,
and let $s$ be a non-CM rational parameter
with $v_p(s)=1$. The other four branch sections are distinct from zero
modulo $p$. The tame local model therefore sends an inertia generator
to $w_{105}$. By the degree character calculation,
\[
 \prod_jm_j^{v_p(a_j)\bmod2}=105
 \quad\hbox{in }\mathbb Q^\times/\mathbb Q^{\times2}.
\]
Here the equality means the norm square class of inertia, and may be
read in the subgroup generated by $2,3,5,7$. All $m_j$ are $p$-units.
The local Hilbert-symbol formula and the fact that $B$ splits at $p$
give
\[
 \operatorname{inv}_p([\operatorname{sign}c])\ne0
 \quad\Longleftrightarrow\quad \left(\frac{105}{p}\right)=-1.
\]
No claim that the sign support is fixed by the discriminant of $B$ is
valid here: the ramification of the degree characters also contributes.

\subsection{A lower bound for every scalar splitting coefficient field}
Suppose a continuous finite-image scalar cochain
$\beta:G_{\mathbb Q}\to\overline{\mathbb Q}^{\times}$, with trivial
action on its coefficient group, satisfies $\delta\beta=c$.
Let $E=\mathbb Q(\beta_\sigma:\sigma\in G_{\mathbb Q})$. Then
\[
 \epsilon_\sigma=\beta_\sigma^2/d_\sigma
\]
is a finite-order character with values in $E$. Moreover
$\beta_\sigma/\sqrt{d_\sigma}$ is a square-root cochain for
$\epsilon$, and its coboundary is $\operatorname{sign}c$.

Put $k=v_2(p-1)$. A finite character of $G_{\mathbb Q_p}$ whose
square-root obstruction is nonzero must have on inertia a value of
order divisible by $2^k$. Here is the local reason. Its wild inertia
part has odd order and hence admits a unique square root, and its
unramified part admits a square root by choosing a square root of its
Frobenius value. For the tame part, the relation
$F\tau F^{-1}=\tau^p$ says that character values on tame inertia lie
in $\mu_{p-1}$. Such a value has a square root in that same group
unless its order has full $2$-part $2^k$. Equivalently, one may use
$\mathbb Q_p^\times=p^{\mathbb Z}\times\mu_{p-1}\times(1+p\mathbb Z_p)$
and local reciprocity.

Thus, when $(105/p)=-1$, the field $E$ contains a primitive
$2^k$-th root of unity (extract the $2$-primary part of that character
value). Consequently
\begin{equation}\label{eq:coefficient-field-obstruction}
 [E:\mathbb Q]\ge2^{k-1}.
\end{equation}
This lower bound applies to every scalar splitting of this cocycle,
whether or not extra ramification is allowed in its field of definition.
In particular it is stronger than an obstruction only to an
unramified-outside-the-old-support splitting.
If such a cochain factors through a finite extension group
$\widetilde\Gamma\to\Gamma$ of the original fiber group, then
$\epsilon$ also factors through $\widetilde\Gamma$. Hence
$2^k$ divides $|\widetilde\Gamma|$. Since $|\Gamma|$ divides $576$,
the kernel has $2$-part at least $2^{k-6}$ when $k\ge6$.
Thus no fixed finite central extension of the fiber group resolves
the scalar splitting problem uniformly either.

\subsection{Parameters forcing unbounded costs}
For any $k\ge2$, impose
\[
 p\equiv1+2^k\pmod{2^{k+1}},\qquad p\equiv71\pmod{105}.
\]
The Chinese remainder theorem gives a reduced residue class, so
Dirichlet's theorem supplies infinitely many such primes, outside
any fixed finite set. They have $v_2(p-1)=k$, and quadratic reciprocity
gives $(105/p)=-1$: the residues modulo $3,5,7$ are respectively
$2,1,1$, while $p\equiv1\pmod4$.

Only finitely many rational parameters of this quotient are CM.
Indeed, a CM point $P\in X$ above a rational parameter has degree at
most $16$. Gonz\'alez--Rotger \cite[Theorem 5.8(1)]{GonzalezRotger2006}
give $H_R=K\mathbb Q(P)$ for its imaginary quadratic order $R$ and
ring class field $H_R$. Hence $h(R)=[H_R:K]\le16$.
Heilbronn's class-number growth theorem \cite{Heilbronn1934} gives finitely many possible
imaginary quadratic fields. In each such field the order class-number
formula implies
\[
 h(R_f)\ge h(\mathcal O_K)\varphi(f)/3\longrightarrow\infty,
\]
so there are finitely many possible orders, and each has finitely many
CM points. Thus for each chosen $p$ one can choose an integer
$s=pu$, with $p\nmid u$, outside the finite CM set and the five branch
values. It has the required valuation-one contact at zero.

Taking $k=14$ in \eqref{eq:coefficient-field-obstruction} already gives
$[E:\mathbb Q]\ge8192>4608$; taking $k$ arbitrarily large makes the
cost unbounded. Therefore the proposed universal bound of 4608
(and the alternative 2304 bound) for scalar splitting coefficient
fields cannot hold for this family with its established branch data.

This conclusion is limited to that scalar-coefficient descent strategy.
It is not a counterexample to the desired five-contact height bound,
nor a proof that every other construction of a bounded-dimensional
auxiliary variety fails. A different height argument would have to
avoid this scalar splitting cost, or a restricted-parameter argument
would have to prove extra local conditions excluding these contacts.
